\section{Responsibilities}

\textit{Normative language per \href{https://tools.ietf.org/html/rfc2119}{RFC~2119}.}

\subsection{Purpose}

This section defines the responsibility boundary between the tenant (acting as the
OAuth~2.0 Authorization Server, 
\href{https://datatracker.ietf.org/doc/html/rfc6749#section-1.1}{RFC~6749~\S1.1}) 
and Canvasflow (acting as the OAuth~2.0 Resource Server, 
\href{https://datatracker.ietf.org/doc/html/rfc6749#section-1.1}{RFC~6749~\S1.1}). 
This matrix is the contractual heart of the integration. Any ambiguity about 
which party is responsible for a given function is resolved here.

\subsection{Tenant Responsibilities (Authorization Server)}

The tenant is the Authorization Server in the sense of 
\href{https://datatracker.ietf.org/doc/html/rfc6749#section-1.1}{RFC~6749~\S1.1}. 
The tenant is solely responsible for the following:

\subsubsection{User Authentication}

The tenant MUST perform all user authentication. This includes, but is not limited to:
credential verification (username and password, passkeys, certificates),
multi-factor authentication (MFA), social identity connections, and session initiation
and termination. Canvasflow has no visibility into the tenant's authentication mechanisms
and makes no assumptions about them.

\subsubsection{Authorization Decisions}

The tenant MUST make all authorization decisions. The determination of which users are
entitled to which Canvasflow resources is entirely the tenant's responsibility. These
decisions are encoded by the pre-token hook at login and refresh time and sealed into
the signed JWT. Canvasflow enforces these decisions; it does not make them.

\subsubsection{Entitlement Mapping}

The tenant MUST maintain the mapping between their internal product or entitlement
identifiers and Canvasflow resource identifiers. This mapping table lives inside the
tenant's own system --- in the IdP hook, a directory attribute, a hardcoded table, or a
lookup against the tenant's own database. The form is the tenant's choice; the output
MUST be an array of Canvasflow resource identifiers in \texttt{cf:<type>:<id>} format
(see \hyperref[sec:mandatory-value-format]{Entitlement Mapping \S4.2}).

This mapping MUST be re-evaluated on every login and every token refresh. The hook MUST
NOT copy the entitlement claim from a prior token; it MUST read from the tenant's live
source of truth on each invocation. This ensures that changes to entitlements propagate
within the revocation window ($\leq$ 3600~seconds, governed by the access token
\texttt{exp}).

\pagebreak

\subsubsection{Token Signing and Key Protection}

The tenant MUST sign all access tokens using a signing method that Canvasflow supports
(\hyperref[sec:signing-path-requirements]{\S5.2}). The tenant MUST protect its signing private keys (JWKS path) or shared secrets
(HS256 path). Canvasflow's guarantee that it cannot forge tokens on behalf of the tenant
depends on the tenant maintaining exclusive control of their private signing keys.

\subsubsection{Key Rotation}

\begin{itemize}
  \item \textbf{JWKS path:} The tenant MAY rotate keys at any time by publishing a new
    key pair under a new \texttt{kid} in their JWKS endpoint. The tenant MUST maintain
    the outgoing key in the JWKS for a minimum overlap window of one hour after the new
    key begins signing tokens. No coordination with Canvasflow is required for JWKS key
    rotation.

  \item \textbf{HS256 path:} The tenant MUST notify Canvasflow before rotating a shared
    secret. The tenant MUST provide the new secret via a secure channel and MUST
    maintain the outgoing secret as valid under its existing \texttt{kid} for a minimum
    of one hour after the new secret takes effect.
\end{itemize}

\subsubsection{Session Lifecycle}\label{sec:session-lifecycle}

The tenant is responsible for login, silent refresh, logout, and revocation of user
sessions. If the tenant needs to immediately cut off a user's access, the tenant MUST
revoke the user's refresh token in the IdP. The residual access window is bounded by
the access token \texttt{exp} ($\leq$ 3600~seconds). Canvasflow has no mechanism to
accelerate this revocation --- that is by design (see \hyperref[sec:out-of-scope]{\S1.4}).



\subsubsection{OAuth Client Registration}

The tenant MUST register the Canvasflow SPA as a public OAuth~2.0 client in their IdP
with the exact settings required in \hyperref[sec:tenant-side-setup]{\S9.3} (Step 2.2), including:

\begin{itemize}
  \item The exact callback URI: \texttt{https://www.<tenant-domain>/callback}
  \item Grant type: Authorization Code with PKCE (S256)
  \item Token endpoint authentication method: none (public client --- no client secret)
  \item Access token format: JWT with \texttt{typ: "at+jwt"} in the header
  \item Access token lifetime: $\leq$ 3600~seconds
\end{itemize}

\pagebreak

\subsection{Canvasflow Responsibilities (Resource Server)}\label{sec:canvasflow-responsibilities}

Canvasflow is the Resource Server in the sense of 
\href{https://datatracker.ietf.org/doc/html/rfc6749#section-1.1}{RFC~6749~\S1.1}. 
Canvasflow is responsible for the following:

\subsubsection{Token Validation on Every Request}

Canvasflow MUST validate the cryptographic signature and all required claims of every
access token received at the GraphQL API. 
% The complete, ordered validation sequence is defined in \texttt{spec\_appendix\_validation}. 
No request is authorized without passing the full sequence.

\subsubsection{Content Gating}

Canvasflow MUST gate GraphQL content strictly by the verified entitlement claim. Content
Canvasflow classifies as free MUST be returned to any bearer of a valid access token,
regardless of the entitlement set. Content Canvasflow classifies as premium MUST be
returned only when the matching \texttt{cf:<type>:<id>} value is present in the
effective entitlement set (see \hyperref[sec:claim-name-hierarchy]{Entitlement Mapping \S4.3}). 
Canvasflow MUST NOT apply any additional authorization logic beyond 
what the signed token encodes.

\subsubsection{JWKS Caching}

Canvasflow MUST fetch tenant JWKS documents server-side only and MUST cache them with a
TTL of one hour. An incoming token whose \texttt{kid} is not found in the cache triggers
an immediate JWKS re-fetch. Canvasflow MUST NOT serve stale JWKS keys when a re-fetch
fails; it MUST reject the token. JWKS key material is never exposed to clients.

\subsubsection{HS256 Secret Storage}

Canvasflow MUST store HS256 shared secrets encrypted at the application layer using
AES-256-GCM with a KMS-managed key. Secrets MUST never appear in logs or error
responses. Canvasflow MUST notify the affected tenant within 24~hours of any confirmed
or suspected compromise of a stored HS256 secret.

\subsubsection{Generic Error Responses}

Canvasflow MUST return generic 401 responses on all token validation failures. The
response MUST NOT reveal which validation step failed, which tenant was involved, or
which claim was absent or malformed. Detailed failure reasons are logged server-side and
correlated by \texttt{jti}. This policy is a security requirement; it prevents both
tenant enumeration and targeted token-forgery attacks.

\subsubsection{Resource ID Stability}

Canvasflow MUST NOT change the Canvasflow resource IDs (for example, \texttt{cf:issue:123})
assigned to a tenant after onboarding without prior written notice to the tenant's
designated technical lead. Resource ID stability is essential because these values are
hardcoded into the tenant's IdP hook mapping table.

\subsubsection{No Runtime Calls to the Tenant IdP}

Canvasflow MUST NOT call the tenant IdP at API request time for any purpose other than
a JWKS cache refresh triggered by an unknown \texttt{kid}. The tenant's IdP availability
MUST NOT be a dependency for serving already-issued access tokens.

\subsection{Justification}

The responsibility distribution above was designed to maximize system reliability and
minimize coupling.

\textbf{Stateless Resource Server.} By confining all authorization logic to the signed
JWT --- a decision sealed by the tenant's IdP before the token reaches Canvasflow --- the
Resource Server requires zero network calls in the hot path (JWKS key cache is warm).
This makes the Canvasflow API horizontally scalable without coordination and removes the
tenant's IdP from the availability dependency chain of every API request.

\textbf{Clean trust boundary.} The signed JWT is the sole interface between the two
systems. A tenant that rotates its signing keys, changes its user directory, or migrates
IdP platforms does so without touching the Canvasflow API layer, provided the output
token structure remains conformant with this specification.

\textbf{Bounded revocation window as a trade-off.} Delegating revocation to the tenant
and bounding the exposure window by $\texttt{exp} \leq 3600\,\mathrm{s}$ is a deliberate
design choice that preserves the stateless Resource Server model. This trade-off is
acceptable for content delivery and is documented here --- not hidden --- so that tenants
can evaluate it against their security requirements.